\section{相关工作}

\paragraph{同策略语言模型蒸馏。}
同策略蒸馏（OPD）通过在稠密教师监督下使用自生成的轨迹进行学生模型训练，缓解了曝光偏差 \citep{gu2024minillm,agarwal2024policy}。后续工作探索了偏斜散度 \citep{ko2024distillm}、自适应目标 \citep{jang2026stable,heo2026policy,jin2026entropy} 以及复杂推理场景中的诊断行为 \citep{fu2026revisiting,zhu2026many}。然而，全词表 OPD 在分布式轨迹生成过程中会产生难以承受的通信和内存开销，严重制约其可扩展性。
\vspace{-4mm}
\paragraph{稀疏散度与截断偏差。}
为缓解通信瓶颈，分布式训练通常采用稀疏的教师端 Top-$k$ 信息包 \citep{zhu2026many}。然而，简单地对 Top-$k$ 支持集进行重归一化或截断，会优化一个有偏的代理目标，导致尾部概率质量分配错误。尽管诸如 TA-OPD \citep{huang2026tail} 等粗支持方法引入了一个虚拟 token \citep{anshumann-etal-2025-sparse} 来聚合被忽略的概率质量，它们仍优化一个代理 $(k+1)$ -way 散度，忽略了教师端尾部内部的细粒度几何结构。
\vspace{-4mm}
\paragraph{随机估计量与方差减小。}
一种替代范式构建了随机估计量，以无偏地恢复真实的全词表 rKL 梯度。标准的采样 token 估计量 \citep{gu2024minillm,lu2025onpolicydistillation} 存在极端的蒙特卡罗方差，促使人们提出如 vOPD \citep{oh2026kl} 这样的方差减少基准方法。HETS 是 EMA-PG \citep{zhang2026ema} 提出的一种混合方案，它对前 $k$ 个 head 进行解析计算，仅在尾部进行采样，但仍在尾部命中时引发严重的梯度突增，导致训练动态不稳定。
